
\documentclass{article}
\usepackage{colortbl}
\usepackage{makecell}
\usepackage{multirow}
\usepackage{supertabular}

\begin{document}

\newcounter{utterance}

\centering \large Interaction Transcript for game `clean\_up', experiment `0\_easy\_7obj\_lv', episode 1 with Qwen3.6{-}35B{-}A3B{-}FP8{-}without{-}reasoning.
\vspace{24pt}

{ \footnotesize  \setcounter{utterance}{1}
\setlength{\tabcolsep}{0pt}
\begin{supertabular}{c@{$\;$}|p{.15\linewidth}@{}p{.15\linewidth}p{.15\linewidth}p{.15\linewidth}p{.15\linewidth}p{.15\linewidth}}
    \# & $\;$A & \multicolumn{4}{c}{Game Master} & $\;\:$B\\
    \hline

    \theutterance \stepcounter{utterance}  
    & & \multicolumn{4}{p{0.6\linewidth}}{
        \cellcolor[rgb]{0.9,0.9,0.9}{
            \makecell[{{p{\linewidth}}}]{
                \texttt{\tiny{[P1$\langle$GM]}}
                \texttt{Es esmu jūsu spēles vadītājs, un jūs spēlējat sadarbības spēli ar šādu režģi kā spēles laukumu:} \\
\\ 
\texttt{```} \\
\texttt{ 1234567} \\
\texttt{╔═══════╗} \\
\texttt{║◌◌◌◌◌T◌║ 1} \\
\texttt{║P◌◌◌◌◌◌║ 2} \\
\texttt{║◌◌AD◌┌─╢ 3} \\
\texttt{║◌◌◌┌─┤S║ 4} \\
\texttt{║O◌◌│◌├─╢ 5} \\
\texttt{╟───┼─┘◌║ 6} \\
\texttt{║◌M◌│◌◌◌║ 7} \\
\texttt{╚═══╧═══╝} \\
\\ 
\texttt{```} \\
\\ 
\texttt{* Augšējā malā ir attēlotas x koordinātas, kas pieaug pa labi, un labajā malā y koordinātas, kas pieaug uz leju.} \\
\texttt{* Uz jūsu režģa ir nejauši izvietoti šādi objekti: 'P', 'O', 'T', 'S', 'D', 'A', 'M'.} \\
\\ 
\\ 
\texttt{Otrs spēlētājs redz spēles laukuma variāciju, kurā objekti ir izvietoti citās nejaušās vietās. Jūs nevarat redzēt otra spēlētāja laukumu, un viņš nevar redzēt jūsējo.} \\
\\ 
\\ 
\texttt{**Mērķis:**} \\
\\ 
\texttt{Abiem spēlētājiem ir jāpārvieto objekti uz sava fona tā, lai vienādi objekti nonāktu vienādās koordinātēs. Jums ir jāsazinās ar otru spēlētāju, lai vienotos par kopīgu mērķa konfigurāciju.} \\
\\ 
\\ 
\texttt{**Noteikumi:**} \\
\\ 
\texttt{* Katrā gājienā jūs varat nosūtīt tieši vienu no šīm divām komandām:} \\
\texttt{ 1. `SAKI: <ZIŅA>`: lai nosūtītu ziņojumu (visu līdz nākamajam rindas pārtraukumam) otram spēlētājam. Es to pārsūtīšu jūsu partnerim.} \\
\texttt{2. `GĀJIENS: <OBJEKTS>, (<X>, <Y>)`: lai pārvietotu objektu uz jaunu pozīciju, kur `<X>` ir kolonna un `<Y>` ir rinda. Es jūs informēšu, vai jūsu gājiens bija veiksmīgs vai nē.} \\
\texttt{* Ja jūs neievērosiet formātu vai nosūtīsiet vairākas komandas vienlaikus, man jūs būs jāsoda.} \\
\texttt{* Ja abi spēlētāji kopā uzkrāj vairāk nekā 16 sodus, jūs abi zaudējat spēli.} \\
\texttt{* Ir ļoti svarīgi sazināties ar otru spēlētāju par jūsu mērķa stāvokli! *Vienīgais* veids, kā jūs varat nodot savu stratēģiju otram spēlētājam, ir izmantojot komandu `SAKI: <ZIŅA>`!} \\
\\ 
\\ 
\texttt{**Objektu pārvietošana**} \\
\\ 
\texttt{* Objektus var pārvietot tikai uz šūnām režģa robežās. Mērķa šūnai jābūt tukšai, t.i., tajā jābūt tikai simbolam '◌'.} \\
\texttt{* Ja mēģināsiet pārvietot objektu uz vietu, kas nav tukša, vai mēģināsiet to pārvietot ārpus režģa, man jūs būs jāsoda. Jūs saņemat vēl vienu mēģinājumu.} \\
\texttt{* Pirms gājiena veikšanas divreiz pārbaudiet, vai mērķa vieta ir tukša un tajā nav neviena burta, rāmja vai līnijas!} \\
\\ 
\\ 
\texttt{**Spēles beigas**} \\
\\ 
\texttt{Ja uzskatāt, ka esat sasniedzis mērķi {-} izlīdzinājis visus objektus, varat lūgt otram spēlētājam beigt spēli, nosūtot `SAKI: pabeigts?`. Ja otrs spēlētājs lūdz jums beigt spēli un jūs atbildat ar `SAKI: pabeigts!`, spēle beigsies.} \\
\\ 
\texttt{Abi spēlētāji uzvar, ja spēle beidzas 28 raundos, kur viens raunds ir definēts kā divi spēlētāji, kas katrs nosūta derīgu komandu.} \\
\\ 
\texttt{**Punktu skaitīšana:**} \\
\\ 
\texttt{Jo tuvāk vienādi objekti atrodas abos spēles laukumos, jo vairāk punktu jūs saņemat. Sodi samazina jūsu punktus. Vai varat pārspēt rekordu?} \\
\\ 
\\ 
\texttt{Lūdzu, nosūtiet ziņu otram spēlētājam, lai sāktu spēli!} \\
            }
        }
    }
    & \\ \\

    \theutterance \stepcounter{utterance}  
    & \multicolumn{4}{p{0.6\linewidth}}{
        \cellcolor[rgb]{0.9,0.9,0.9}{
            \makecell[{{p{\linewidth}}}]{
                \texttt{\tiny{[P1$\rangle$GM]}}
                \texttt{SAKI: Sveiki! Lai veiksmīgi saskaņotu objektus, man ir nepieciešams zināt jūsu objektu sākotnējās koordinātas. Lūdzu, nosūtiet man sava laukuma objektu 'P', 'O', 'T', 'S', 'D', 'A', 'M' pašreizējās (x, y) koordinātas. Es nosūtīšu savas koordinātas pēc tam, kad saņemšu jūsu datus.} \\
            }
        }
    }
    & & \\ \\

    \theutterance \stepcounter{utterance}  
    & & & \multicolumn{2}{p{0.3\linewidth}}{
        \cellcolor[rgb]{0.9,0.9,0.9}{
            \makecell[{{p{\linewidth}}}]{
                \texttt{\tiny{[GM$|$GM]}}
                \texttt{SAKI: Sveiki! Lai veiksmīgi saskaņotu objektus, man ir nepieciešams zināt jūsu objektu sākotnējās koordinātas. Lūdzu, nosūtiet man sava laukuma objektu 'P', 'O', 'T', 'S', 'D', 'A', 'M' pašreizējās (x, y) koordinātas. Es nosūtīšu savas koordinātas pēc tam, kad saņemšu jūsu datus.} \\
            }
        }
    }
    & & \\ \\

    \theutterance \stepcounter{utterance}  
    & & \multicolumn{4}{p{0.6\linewidth}}{
        \cellcolor[rgb]{0.9,0.9,0.9}{
            \makecell[{{p{\linewidth}}}]{
                \texttt{\tiny{[GM$\rangle$P2]}}
                \texttt{Es esmu jūsu spēles vadītājs, un jūs spēlējat sadarbības spēli ar šādu režģi kā spēles laukumu:} \\
\\ 
\texttt{```} \\
\texttt{ 1234567} \\
\texttt{╔═══════╗} \\
\texttt{║◌◌◌◌◌S◌║ 1} \\
\texttt{║◌◌◌◌O◌◌║ 2} \\
\texttt{║◌◌◌M◌┌─╢ 3} \\
\texttt{║◌◌A┌─┤◌║ 4} \\
\texttt{║◌◌◌│◌├─╢ 5} \\
\texttt{╟───┼─┘◌║ 6} \\
\texttt{║D◌T│P◌◌║ 7} \\
\texttt{╚═══╧═══╝} \\
\\ 
\texttt{```} \\
\\ 
\texttt{* Augšējā malā ir attēlotas x koordinātas, kas pieaug pa labi, un labajā malā y koordinātas, kas pieaug uz leju.} \\
\texttt{* Uz jūsu režģa ir nejauši izvietoti šādi objekti: 'P', 'O', 'T', 'S', 'D', 'A', 'M'.} \\
\\ 
\\ 
\texttt{Otrs spēlētājs redz spēles laukuma variāciju, kurā objekti ir izvietoti citās nejaušās vietās. Jūs nevarat redzēt otra spēlētāja laukumu, un viņš nevar redzēt jūsējo.} \\
\\ 
\\ 
\texttt{**Mērķis:**} \\
\\ 
\texttt{Abiem spēlētājiem ir jāpārvieto objekti uz sava fona tā, lai vienādi objekti nonāktu vienādās koordinātēs. Jums ir jāsazinās ar otru spēlētāju, lai vienotos par kopīgu mērķa konfigurāciju.} \\
\\ 
\\ 
\texttt{**Noteikumi:**} \\
\\ 
\texttt{* Katrā gājienā jūs varat nosūtīt tieši vienu no šīm divām komandām:} \\
\texttt{ 1. `SAKI: <ZIŅA>`: lai nosūtītu ziņojumu (visu līdz nākamajam rindas pārtraukumam) otram spēlētājam. Es to pārsūtīšu jūsu partnerim.} \\
\texttt{2. `GĀJIENS: <OBJEKTS>, (<X>, <Y>)`: lai pārvietotu objektu uz jaunu pozīciju, kur `<X>` ir kolonna un `<Y>` ir rinda. Es jūs informēšu, vai jūsu gājiens bija veiksmīgs vai nē.} \\
\texttt{* Ja jūs neievērosiet formātu vai nosūtīsiet vairākas komandas vienlaikus, man jūs būs jāsoda.} \\
\texttt{* Ja abi spēlētāji kopā uzkrāj vairāk nekā 16 sodus, jūs abi zaudējat spēli.} \\
\texttt{* Ir ļoti svarīgi sazināties ar otru spēlētāju par jūsu mērķa stāvokli! *Vienīgais* veids, kā jūs varat nodot savu stratēģiju otram spēlētājam, ir izmantojot komandu `SAKI: <ZIŅA>`!} \\
\\ 
\\ 
\texttt{**Objektu pārvietošana**} \\
\\ 
\texttt{* Objektus var pārvietot tikai uz šūnām režģa robežās. Mērķa šūnai jābūt tukšai, t.i., tajā jābūt tikai simbolam '◌'.} \\
\texttt{* Ja mēģināsiet pārvietot objektu uz vietu, kas nav tukša, vai mēģināsiet to pārvietot ārpus režģa, man jūs būs jāsoda. Jūs saņemat vēl vienu mēģinājumu.} \\
\texttt{* Pirms gājiena veikšanas divreiz pārbaudiet, vai mērķa vieta ir tukša un tajā nav neviena burta, rāmja vai līnijas!} \\
\\ 
\\ 
\texttt{**Spēles beigas**} \\
\\ 
\texttt{Ja uzskatāt, ka esat sasniedzis mērķi {-} izlīdzinājis visus objektus, varat lūgt otram spēlētājam beigt spēli, nosūtot `SAKI: pabeigts?`. Ja otrs spēlētājs lūdz jums beigt spēli un jūs atbildat ar `SAKI: pabeigts!`, spēle beigsies.} \\
\\ 
\texttt{Abi spēlētāji uzvar, ja spēle beidzas 28 raundos, kur viens raunds ir definēts kā divi spēlētāji, kas katrs nosūta derīgu komandu.} \\
\\ 
\texttt{**Punktu skaitīšana:**} \\
\\ 
\texttt{Jo tuvāk vienādi objekti atrodas abos spēles laukumos, jo vairāk punktu jūs saņemat. Sodi samazina jūsu punktus. Vai varat pārspēt rekordu?} \\
\\ 
\texttt{Otrs spēlētājs sāka spēli, nosūtot šo ziņojumu:} \\
\\ 
\texttt{"Sveiki! Lai veiksmīgi saskaņotu objektus, man ir nepieciešams zināt jūsu objektu sākotnējās koordinātas. Lūdzu, nosūtiet man sava laukuma objektu 'P', 'O', 'T', 'S', 'D', 'A', 'M' pašreizējās (x, y) koordinātas. Es nosūtīšu savas koordinātas pēc tam, kad saņemšu jūsu datus."} \\
\\ 
\texttt{Kāda ir jūsu pirmā komanda?} \\
            }
        }
    }
    & \\ \\

    \theutterance \stepcounter{utterance}  
    & & & \multicolumn{4}{p{0.6\linewidth}}{
        \cellcolor[rgb]{0.9,0.9,0.9}{
            \makecell[{{p{\linewidth}}}]{
                \texttt{\tiny{[GM$\langle$P2]}}
                \texttt{SAKI: Sveiki! Lūk, manu objektu koordinātas: P(6,7), O(5,2), T(4,7), S(6,1), D(1,7), A(3,4), M(4,3). Lūdzu, nosūtiet savas koordinātas, lai mēs varam vienoties par mērķa konfigurāciju.} \\
            }
        }
    }
     \\ \\

    \theutterance \stepcounter{utterance}  
    & & & \multicolumn{2}{p{0.3\linewidth}}{
        \cellcolor[rgb]{0.9,0.9,0.9}{
            \makecell[{{p{\linewidth}}}]{
                \texttt{\tiny{[GM$|$GM]}}
                \texttt{SAKI: Sveiki! Lūk, manu objektu koordinātas: P(6,7), O(5,2), T(4,7), S(6,1), D(1,7), A(3,4), M(4,3). Lūdzu, nosūtiet savas koordinātas, lai mēs varam vienoties par mērķa konfigurāciju.} \\
            }
        }
    }
    & & \\ \\

    \theutterance \stepcounter{utterance}  
    & & \multicolumn{4}{p{0.6\linewidth}}{
        \cellcolor[rgb]{0.9,0.9,0.9}{
            \makecell[{{p{\linewidth}}}]{
                \texttt{\tiny{[P1$\langle$GM]}}
                \texttt{Jūsu ziņojums ir nodots otram spēlētājam.} \\
\\ 
\\ 
\texttt{Pašlaik jūs spēlējat kārtu 1 no maksimāli 28.} \\
\texttt{Jūs kopā esat uzkrājuši 0 no 16 sodiem.} \\
\texttt{Otrs spēlētājs nosūtīja šo ziņojumu:} \\
\\ 
\\ 
\texttt{"Sveiki! Lūk, manu objektu koordinātas: P(6,7), O(5,2), T(4,7), S(6,1), D(1,7), A(3,4), M(4,3). Lūdzu, nosūtiet savas koordinātas, lai mēs varam vienoties par mērķa konfigurāciju."} \\
\\ 
\\ 
\texttt{Kāda ir jūsu nākamā komanda?} \\
            }
        }
    }
    & \\ \\

    \theutterance \stepcounter{utterance}  
    & \multicolumn{4}{p{0.6\linewidth}}{
        \cellcolor[rgb]{0.9,0.9,0.9}{
            \makecell[{{p{\linewidth}}}]{
                \texttt{\tiny{[P1$\rangle$GM]}}
                \texttt{SAKI: Paldies! Lūk, manas objektu koordinātas: P(1,2), O(1,5), T(6,1), S(6,4), D(3,3), A(3,3), M(2,7).} \\
\\ 
\texttt{Lai vienkāršotu saskaņošanu, ierosinu mērķa konfigurāciju, kurā katrs objekts tiek pārvietots uz tā pašas tipa objekta (pēc alfabēta vai kārtas) vismazāko attālumu, vai vienkārši fiksēsim vienu kopīgu shēmu.} \\
\\ 
\texttt{Ņemot vērā, ka mums ir divi 'A' un 'D' objekti (manējie A un D atrodas vienā šūnā (3,3), kas ir kļūda? Nē, skatāmies rūgti:} \\
\texttt{Manā režģī:} \\
\texttt{Rinda 3: `◌◌AD◌┌─╢` {-}> A ir (3,3), D ir (4,3)? Nē, skaitot kolonnas:} \\
\texttt{1: ◌} \\
\texttt{2: ◌} \\
\texttt{3: A} \\
\texttt{4: D} \\
\texttt{Tātad manējie: A(3,3), D(4,3).} \\
\\ 
\texttt{Pārbaudīsim manu sarakstu no iepriekšējā ziņojuma, ko es "izdomāju" kā atbildi, bet man jābūt precīzam.} \\
\texttt{Manā laukumā:} \\
\texttt{P: (1,2)} \\
\texttt{O: (1,5)} \\
\texttt{T: (6,1)} \\
\texttt{S: (6,4)} \\
\texttt{A: (3,3)} \\
\texttt{D: (4,3)} \\
\texttt{M: (2,7)} \\
\\ 
\texttt{Otra spēlētāja laukums:} \\
\texttt{P: (6,7)} \\
\texttt{O: (5,2)} \\
\texttt{T: (4,7)} \\
\texttt{S: (6,1)} \\
\texttt{D: (1,7)} \\
\texttt{A: (3,4)} \\
\texttt{M: (4,3)} \\
\\ 
\texttt{Lai vienkāršotu, ierosinu mērķa pozīcijas, kurās mēs abus objektus (no abiem spēlētājiem) novietosim uz vienu un to pašu koordinātu. Tā kā mēs nevaram redzēt viens otra laukumu, mums ir jāvienojas par *koordinātām*, nevis par to, kurš objektu pārvieto. Bet uzdevums saka: "Abiem spēlētājiem ir jāpārvieto objekti uz sava fona tā, lai vienādi objekti nonāktu vienādās koordinātēs." Tas nozīmē, ka manam 'P' un partnera 'P' ir jābūt vienādās (x,y) koordinātēs.} \\
\\ 
\texttt{Tā kā mēs nevaram redzēt, vai mērķa šūna ir tukša otram spēlētājam, mums ir jāizvēlas tādas koordinātas, kas ir drošas vai vienkārši jāvienojas par "mērķa punktiem". Bet uzdevuma noteikumi saka: "Mērķa šūnai jābūt tukšai". Tas attiecas uz *jūsu* laukumu. Partnerim tas attiecas uz *viņa* laukumu.} \\
\\ 
\texttt{Tā kā mēs nevaram zināt, kuras šūnas ir tukšas otram, mums ir jāizvēlas tādas koordinātas, kas ir mazāk aizņētas, vai vienkārši jāvienojas par konkrētām koordinātām un cerēt, ka tās ir pieejamas.} \\
\\ 
\texttt{Ierosinu šādu mērķa konfigurāciju (visiem objektiem):} \\
\texttt{P {-}> (1,1)} \\
\texttt{O {-}> (1,2)} \\
\texttt{T {-}> (1,3)} \\
\texttt{S {-}> (1,4)} \\
\texttt{D {-}> (1,5)} \\
\texttt{A {-}> (1,6)} \\
\texttt{M {-}> (1,7)} \\
\\ 
\texttt{Vai arī vienkāršāk: katrs objekts iet uz savu "sākuma" pozīciju, ja tā ir tukša? Nē, mums jāvienojas.} \\
\\ 
\texttt{Ierosinu:} \\
\texttt{P: (1,1)} \\
\texttt{O: (2,1)} \\
\texttt{T: (3,1)} \\
\texttt{S: (4,1)} \\
\texttt{D: (5,1)} \\
\texttt{A: (6,1)} \\
\texttt{M: (1,2)} \\
\\ 
\texttt{Šī ir tikai viena iespēja. Lūdzu, apstipriniet, vai jūs piekrītat šai mērķa konfigurācijai, vai arī ierosiniet savu.} \\
\\ 
\texttt{SAKI: Apstipriniet šo mērķa konfigurāciju: P(1,1), O(2,1), T(3,1), S(4,1), D(5,1), A(6,1), M(1,2). Ja piekrītat, es sāksim pārvietot objektus.} \\
            }
        }
    }
    & & \\ \\

    \theutterance \stepcounter{utterance}  
    & & & \multicolumn{2}{p{0.3\linewidth}}{
        \cellcolor[rgb]{0.9,0.9,0.9}{
            \makecell[{{p{\linewidth}}}]{
                \texttt{\tiny{[GM$|$GM]}}
                \texttt{SAKI: Paldies! Lūk, manas objektu koordinātas: P(1,2), O(1,5), T(6,1), S(6,4), D(3,3), A(3,3), M(2,7).} \\
\\ 
\texttt{Lai vienkāršotu saskaņošanu, ierosinu mērķa konfigurāciju, kurā katrs objekts tiek pārvietots uz tā pašas tipa objekta (pēc alfabēta vai kārtas) vismazāko attālumu, vai vienkārši fiksēsim vienu kopīgu shēmu.} \\
\\ 
\texttt{Ņemot vērā, ka mums ir divi 'A' un 'D' objekti (manējie A un D atrodas vienā šūnā (3,3), kas ir kļūda? Nē, skatāmies rūgti:} \\
\texttt{Manā režģī:} \\
\texttt{Rinda 3: `◌◌AD◌┌─╢` {-}> A ir (3,3), D ir (4,3)? Nē, skaitot kolonnas:} \\
\texttt{1: ◌} \\
\texttt{2: ◌} \\
\texttt{3: A} \\
\texttt{4: D} \\
\texttt{Tātad manējie: A(3,3), D(4,3).} \\
\\ 
\texttt{Pārbaudīsim manu sarakstu no iepriekšējā ziņojuma, ko es "izdomāju" kā atbildi, bet man jābūt precīzam.} \\
\texttt{Manā laukumā:} \\
\texttt{P: (1,2)} \\
\texttt{O: (1,5)} \\
\texttt{T: (6,1)} \\
\texttt{S: (6,4)} \\
\texttt{A: (3,3)} \\
\texttt{D: (4,3)} \\
\texttt{M: (2,7)} \\
\\ 
\texttt{Otra spēlētāja laukums:} \\
\texttt{P: (6,7)} \\
\texttt{O: (5,2)} \\
\texttt{T: (4,7)} \\
\texttt{S: (6,1)} \\
\texttt{D: (1,7)} \\
\texttt{A: (3,4)} \\
\texttt{M: (4,3)} \\
\\ 
\texttt{Lai vienkāršotu, ierosinu mērķa pozīcijas, kurās mēs abus objektus (no abiem spēlētājiem) novietosim uz vienu un to pašu koordinātu. Tā kā mēs nevaram redzēt viens otra laukumu, mums ir jāvienojas par *koordinātām*, nevis par to, kurš objektu pārvieto. Bet uzdevums saka: "Abiem spēlētājiem ir jāpārvieto objekti uz sava fona tā, lai vienādi objekti nonāktu vienādās koordinātēs." Tas nozīmē, ka manam 'P' un partnera 'P' ir jābūt vienādās (x,y) koordinātēs.} \\
\\ 
\texttt{Tā kā mēs nevaram redzēt, vai mērķa šūna ir tukša otram spēlētājam, mums ir jāizvēlas tādas koordinātas, kas ir drošas vai vienkārši jāvienojas par "mērķa punktiem". Bet uzdevuma noteikumi saka: "Mērķa šūnai jābūt tukšai". Tas attiecas uz *jūsu* laukumu. Partnerim tas attiecas uz *viņa* laukumu.} \\
\\ 
\texttt{Tā kā mēs nevaram zināt, kuras šūnas ir tukšas otram, mums ir jāizvēlas tādas koordinātas, kas ir mazāk aizņētas, vai vienkārši jāvienojas par konkrētām koordinātām un cerēt, ka tās ir pieejamas.} \\
\\ 
\texttt{Ierosinu šādu mērķa konfigurāciju (visiem objektiem):} \\
\texttt{P {-}> (1,1)} \\
\texttt{O {-}> (1,2)} \\
\texttt{T {-}> (1,3)} \\
\texttt{S {-}> (1,4)} \\
\texttt{D {-}> (1,5)} \\
\texttt{A {-}> (1,6)} \\
\texttt{M {-}> (1,7)} \\
\\ 
\texttt{Vai arī vienkāršāk: katrs objekts iet uz savu "sākuma" pozīciju, ja tā ir tukša? Nē, mums jāvienojas.} \\
\\ 
\texttt{Ierosinu:} \\
\texttt{P: (1,1)} \\
\texttt{O: (2,1)} \\
\texttt{T: (3,1)} \\
\texttt{S: (4,1)} \\
\texttt{D: (5,1)} \\
\texttt{A: (6,1)} \\
\texttt{M: (1,2)} \\
\\ 
\texttt{Šī ir tikai viena iespēja. Lūdzu, apstipriniet, vai jūs piekrītat šai mērķa konfigurācijai, vai arī ierosiniet savu.} \\
\\ 
\texttt{SAKI: Apstipriniet šo mērķa konfigurāciju: P(1,1), O(2,1), T(3,1), S(4,1), D(5,1), A(6,1), M(1,2). Ja piekrītat, es sāksim pārvietot objektus.} \\
            }
        }
    }
    & & \\ \\

    \theutterance \stepcounter{utterance}  
    & & & \multicolumn{2}{p{0.3\linewidth}}{
        \cellcolor[rgb]{0.9,0.9,0.9}{
            \makecell[{{p{\linewidth}}}]{
                \texttt{\tiny{[GM$|$GM]}}
                \texttt{Invalid format: head is not empty: 'SAKI: Paldies! Lūk, manas objektu koordinātas: P(1,2), O(1,5), T(6,1), S(6,4), D(3,3), A(3,3), M(2,7).} \\
\\ 
\texttt{Lai vienkāršotu saskaņošanu, ierosinu mērķa konfigurāciju, kurā katrs objekts tiek pārvietots uz tā pašas tipa objekta (pēc alfabēta vai kārtas) vismazāko attālumu, vai vienkārši fiksēsim vienu kopīgu shēmu.} \\
\\ 
\texttt{Ņemot vērā, ka mums ir divi 'A' un 'D' objekti (manējie A un D atrodas vienā šūnā (3,3), kas ir kļūda? Nē, skatāmies rūgti:} \\
\texttt{Manā režģī:} \\
\texttt{Rinda 3: ◌◌AD◌┌─╢ {-}> A ir (3,3), D ir (4,3)? Nē, skaitot kolonnas:} \\
\texttt{1: ◌} \\
\texttt{2: ◌} \\
\texttt{3: A} \\
\texttt{4: D} \\
\texttt{Tātad manējie: A(3,3), D(4,3).} \\
\\ 
\texttt{Pārbaudīsim manu sarakstu no iepriekšējā ziņojuma, ko es "izdomāju" kā atbildi, bet man jābūt precīzam.} \\
\texttt{Manā laukumā:} \\
\texttt{P: (1,2)} \\
\texttt{O: (1,5)} \\
\texttt{T: (6,1)} \\
\texttt{S: (6,4)} \\
\texttt{A: (3,3)} \\
\texttt{D: (4,3)} \\
\texttt{M: (2,7)} \\
\\ 
\texttt{Otra spēlētāja laukums:} \\
\texttt{P: (6,7)} \\
\texttt{O: (5,2)} \\
\texttt{T: (4,7)} \\
\texttt{S: (6,1)} \\
\texttt{D: (1,7)} \\
\texttt{A: (3,4)} \\
\texttt{M: (4,3)} \\
\\ 
\texttt{Lai vienkāršotu, ierosinu mērķa pozīcijas, kurās mēs abus objektus (no abiem spēlētājiem) novietosim uz vienu un to pašu koordinātu. Tā kā mēs nevaram redzēt viens otra laukumu, mums ir jāvienojas par *koordinātām*, nevis par to, kurš objektu pārvieto. Bet uzdevums saka: "Abiem spēlētājiem ir jāpārvieto objekti uz sava fona tā, lai vienādi objekti nonāktu vienādās koordinātēs." Tas nozīmē, ka manam 'P' un partnera 'P' ir jābūt vienādās (x,y) koordinātēs.} \\
\\ 
\texttt{Tā kā mēs nevaram redzēt, vai mērķa šūna ir tukša otram spēlētājam, mums ir jāizvēlas tādas koordinātas, kas ir drošas vai vienkārši jāvienojas par "mērķa punktiem". Bet uzdevuma noteikumi saka: "Mērķa šūnai jābūt tukšai". Tas attiecas uz *jūsu* laukumu. Partnerim tas attiecas uz *viņa* laukumu.} \\
\\ 
\texttt{Tā kā mēs nevaram zināt, kuras šūnas ir tukšas otram, mums ir jāizvēlas tādas koordinātas, kas ir mazāk aizņētas, vai vienkārši jāvienojas par konkrētām koordinātām un cerēt, ka tās ir pieejamas.} \\
\\ 
\texttt{Ierosinu šādu mērķa konfigurāciju (visiem objektiem):} \\
\texttt{P {-}> (1,1)} \\
\texttt{O {-}> (1,2)} \\
\texttt{T {-}> (1,3)} \\
\texttt{S {-}> (1,4)} \\
\texttt{D {-}> (1,5)} \\
\texttt{A {-}> (1,6)} \\
\texttt{M {-}> (1,7)} \\
\\ 
\texttt{Vai arī vienkāršāk: katrs objekts iet uz savu "sākuma" pozīciju, ja tā ir tukša? Nē, mums jāvienojas.} \\
\\ 
\texttt{Ierosinu:} \\
\texttt{P: (1,1)} \\
\texttt{O: (2,1)} \\
\texttt{T: (3,1)} \\
\texttt{S: (4,1)} \\
\texttt{D: (5,1)} \\
\texttt{A: (6,1)} \\
\texttt{M: (1,2)} \\
\\ 
\texttt{Šī ir tikai viena iespēja. Lūdzu, apstipriniet, vai jūs piekrītat šai mērķa konfigurācijai, vai arī ierosiniet savu.} \\
\\ 
\texttt{'} \\
            }
        }
    }
    & & \\ \\

    \theutterance \stepcounter{utterance}  
    & & \multicolumn{4}{p{0.6\linewidth}}{
        \cellcolor[rgb]{0.9,0.9,0.9}{
            \makecell[{{p{\linewidth}}}]{
                \texttt{\tiny{[P1$\langle$GM]}}
                \texttt{Sods: Jūsu ziņojumā nedrīkst būt nekas pirms komandas!} \\
\\ 
\texttt{Pārliecinieties, ka jūsu atbildē ir tikai `SAKI: <ZIŅA>` vai `GĀJIENS: <OBJEKTS>, (<X>, <Y>)`, un nekas cits!} \\
\\ 
\texttt{Jūs kopā esat uzkrājuši 1 no 16 sodiem.} \\
\texttt{Lūdzu, mēģiniet vēlreiz!} \\
            }
        }
    }
    & \\ \\

    \theutterance \stepcounter{utterance}  
    & \multicolumn{4}{p{0.6\linewidth}}{
        \cellcolor[rgb]{0.9,0.9,0.9}{
            \makecell[{{p{\linewidth}}}]{
                \texttt{\tiny{[P1$\rangle$GM]}}
                \texttt{SAKI: Manas koordinātas: P(1,2), O(1,5), T(6,1), S(6,4), A(3,3), D(4,3), M(2,7). Ierosinu mērķa pozīcijas: P(1,1), O(2,1), T(3,1), S(4,1), D(5,1), A(6,1), M(1,2). Vai piekrītat?} \\
            }
        }
    }
    & & \\ \\

    \theutterance \stepcounter{utterance}  
    & & & \multicolumn{2}{p{0.3\linewidth}}{
        \cellcolor[rgb]{0.9,0.9,0.9}{
            \makecell[{{p{\linewidth}}}]{
                \texttt{\tiny{[GM$|$GM]}}
                \texttt{SAKI: Manas koordinātas: P(1,2), O(1,5), T(6,1), S(6,4), A(3,3), D(4,3), M(2,7). Ierosinu mērķa pozīcijas: P(1,1), O(2,1), T(3,1), S(4,1), D(5,1), A(6,1), M(1,2). Vai piekrītat?} \\
            }
        }
    }
    & & \\ \\

    \theutterance \stepcounter{utterance}  
    & & \multicolumn{4}{p{0.6\linewidth}}{
        \cellcolor[rgb]{0.9,0.9,0.9}{
            \makecell[{{p{\linewidth}}}]{
                \texttt{\tiny{[GM$\rangle$P2]}}
                \texttt{Jūsu ziņojums ir nodots otram spēlētājam.} \\
\\ 
\\ 
\texttt{Pašlaik jūs spēlējat kārtu 2 no maksimāli 28.} \\
\texttt{Jūs kopā esat uzkrājuši 1 no 16 sodiem.} \\
\texttt{Otrs spēlētājs nosūtīja šo ziņojumu:} \\
\\ 
\\ 
\texttt{"Manas koordinātas: P(1,2), O(1,5), T(6,1), S(6,4), A(3,3), D(4,3), M(2,7). Ierosinu mērķa pozīcijas: P(1,1), O(2,1), T(3,1), S(4,1), D(5,1), A(6,1), M(1,2). Vai piekrītat?"} \\
\\ 
\\ 
\texttt{Kāda ir jūsu nākamā komanda?} \\
            }
        }
    }
    & \\ \\

    \theutterance \stepcounter{utterance}  
    & & & \multicolumn{4}{p{0.6\linewidth}}{
        \cellcolor[rgb]{0.9,0.9,0.9}{
            \makecell[{{p{\linewidth}}}]{
                \texttt{\tiny{[GM$\langle$P2]}}
                \texttt{SAKI: Jā, es piekrītu šai mērķa konfigurācijai. Tagad sāksim pārvietot objektus. Es sāku ar objektu 'S', lai novietotu to uz (4,1).} \\
            }
        }
    }
     \\ \\

    \theutterance \stepcounter{utterance}  
    & & & \multicolumn{2}{p{0.3\linewidth}}{
        \cellcolor[rgb]{0.9,0.9,0.9}{
            \makecell[{{p{\linewidth}}}]{
                \texttt{\tiny{[GM$|$GM]}}
                \texttt{SAKI: Jā, es piekrītu šai mērķa konfigurācijai. Tagad sāksim pārvietot objektus. Es sāku ar objektu 'S', lai novietotu to uz (4,1).} \\
            }
        }
    }
    & & \\ \\

    \theutterance \stepcounter{utterance}  
    & & \multicolumn{4}{p{0.6\linewidth}}{
        \cellcolor[rgb]{0.9,0.9,0.9}{
            \makecell[{{p{\linewidth}}}]{
                \texttt{\tiny{[P1$\langle$GM]}}
                \texttt{Jūsu ziņojums ir nodots otram spēlētājam.} \\
\\ 
\\ 
\texttt{Pašlaik jūs spēlējat kārtu 2 no maksimāli 28.} \\
\texttt{Jūs kopā esat uzkrājuši 1 no 16 sodiem.} \\
\texttt{Otrs spēlētājs nosūtīja šo ziņojumu:} \\
\\ 
\\ 
\texttt{"Jā, es piekrītu šai mērķa konfigurācijai. Tagad sāksim pārvietot objektus. Es sāku ar objektu 'S', lai novietotu to uz (4,1)."} \\
\\ 
\\ 
\texttt{Kāda ir jūsu nākamā komanda?} \\
            }
        }
    }
    & \\ \\

    \theutterance \stepcounter{utterance}  
    & \multicolumn{4}{p{0.6\linewidth}}{
        \cellcolor[rgb]{0.9,0.9,0.9}{
            \makecell[{{p{\linewidth}}}]{
                \texttt{\tiny{[P1$\rangle$GM]}}
                \texttt{GĀJIENS: S, (4,1)} \\
            }
        }
    }
    & & \\ \\

    \theutterance \stepcounter{utterance}  
    & & & \multicolumn{2}{p{0.3\linewidth}}{
        \cellcolor[rgb]{0.9,0.9,0.9}{
            \makecell[{{p{\linewidth}}}]{
                \texttt{\tiny{[GM$|$GM]}}
                \texttt{GĀJIENS: S, (4,1)} \\
            }
        }
    }
    & & \\ \\

    \theutterance \stepcounter{utterance}  
    & & & \multicolumn{2}{p{0.3\linewidth}}{
        \cellcolor[rgb]{0.9,0.9,0.9}{
            \makecell[{{p{\linewidth}}}]{
                \texttt{\tiny{[GM$|$GM]}}
                \texttt{S moved to (4, 1)} \\
            }
        }
    }
    & & \\ \\

    \theutterance \stepcounter{utterance}  
    & & \multicolumn{4}{p{0.6\linewidth}}{
        \cellcolor[rgb]{0.9,0.9,0.9}{
            \makecell[{{p{\linewidth}}}]{
                \texttt{\tiny{[GM$\rangle$P2]}}
                \texttt{Jūsu ziņojums ir nodots otram spēlētājam.} \\
\\ 
\\ 
\texttt{Pašlaik jūs spēlējat kārtu 3 no maksimāli 28.} \\
\texttt{Jūs kopā esat uzkrājuši 1 no 16 sodiem.} \\
\texttt{Otrs spēlētājs pārvietoja objektu savā režģī.} \\
\\ 
\\ 
\texttt{Kāda ir jūsu nākamā komanda?} \\
            }
        }
    }
    & \\ \\

    \theutterance \stepcounter{utterance}  
    & & & \multicolumn{4}{p{0.6\linewidth}}{
        \cellcolor[rgb]{0.9,0.9,0.9}{
            \makecell[{{p{\linewidth}}}]{
                \texttt{\tiny{[GM$\langle$P2]}}
                \texttt{GĀJIENS: S, (4,1)} \\
            }
        }
    }
     \\ \\

    \theutterance \stepcounter{utterance}  
    & & & \multicolumn{2}{p{0.3\linewidth}}{
        \cellcolor[rgb]{0.9,0.9,0.9}{
            \makecell[{{p{\linewidth}}}]{
                \texttt{\tiny{[GM$|$GM]}}
                \texttt{GĀJIENS: S, (4,1)} \\
            }
        }
    }
    & & \\ \\

    \theutterance \stepcounter{utterance}  
    & & & \multicolumn{2}{p{0.3\linewidth}}{
        \cellcolor[rgb]{0.9,0.9,0.9}{
            \makecell[{{p{\linewidth}}}]{
                \texttt{\tiny{[GM$|$GM]}}
                \texttt{S moved to (4, 1)} \\
            }
        }
    }
    & & \\ \\

    \theutterance \stepcounter{utterance}  
    & & \multicolumn{4}{p{0.6\linewidth}}{
        \cellcolor[rgb]{0.9,0.9,0.9}{
            \makecell[{{p{\linewidth}}}]{
                \texttt{\tiny{[P1$\langle$GM]}}
                \texttt{Veiksmīgi pārvietots 'S' uz (4,1). Jūsu atjauninātais režģis izskatās šādi:} \\
\texttt{```} \\
\texttt{ 1234567} \\
\texttt{╔═══════╗} \\
\texttt{║◌◌◌S◌T◌║ 1} \\
\texttt{║P◌◌◌◌◌◌║ 2} \\
\texttt{║◌◌AD◌┌─╢ 3} \\
\texttt{║◌◌◌┌─┤◌║ 4} \\
\texttt{║O◌◌│◌├─╢ 5} \\
\texttt{╟───┼─┘◌║ 6} \\
\texttt{║◌M◌│◌◌◌║ 7} \\
\texttt{╚═══╧═══╝} \\
\texttt{```} \\
\texttt{Pašlaik jūs spēlējat kārtu 3 no maksimāli 28.} \\
\texttt{Jūs kopā esat uzkrājuši 1 no 16 sodiem.} \\
\texttt{Otrs spēlētājs pārvietoja objektu savā režģī.} \\
\\ 
\\ 
\texttt{Kāda ir jūsu nākamā komanda?} \\
            }
        }
    }
    & \\ \\

    \theutterance \stepcounter{utterance}  
    & \multicolumn{4}{p{0.6\linewidth}}{
        \cellcolor[rgb]{0.9,0.9,0.9}{
            \makecell[{{p{\linewidth}}}]{
                \texttt{\tiny{[P1$\rangle$GM]}}
                \texttt{GĀJIENS: P, (1,1)} \\
            }
        }
    }
    & & \\ \\

    \theutterance \stepcounter{utterance}  
    & & & \multicolumn{2}{p{0.3\linewidth}}{
        \cellcolor[rgb]{0.9,0.9,0.9}{
            \makecell[{{p{\linewidth}}}]{
                \texttt{\tiny{[GM$|$GM]}}
                \texttt{GĀJIENS: P, (1,1)} \\
            }
        }
    }
    & & \\ \\

    \theutterance \stepcounter{utterance}  
    & & & \multicolumn{2}{p{0.3\linewidth}}{
        \cellcolor[rgb]{0.9,0.9,0.9}{
            \makecell[{{p{\linewidth}}}]{
                \texttt{\tiny{[GM$|$GM]}}
                \texttt{P moved to (1, 1)} \\
            }
        }
    }
    & & \\ \\

    \theutterance \stepcounter{utterance}  
    & & \multicolumn{4}{p{0.6\linewidth}}{
        \cellcolor[rgb]{0.9,0.9,0.9}{
            \makecell[{{p{\linewidth}}}]{
                \texttt{\tiny{[GM$\rangle$P2]}}
                \texttt{Veiksmīgi pārvietots 'S' uz (4,1). Jūsu atjauninātais režģis izskatās šādi:} \\
\texttt{```} \\
\texttt{ 1234567} \\
\texttt{╔═══════╗} \\
\texttt{║◌◌◌S◌◌◌║ 1} \\
\texttt{║◌◌◌◌O◌◌║ 2} \\
\texttt{║◌◌◌M◌┌─╢ 3} \\
\texttt{║◌◌A┌─┤◌║ 4} \\
\texttt{║◌◌◌│◌├─╢ 5} \\
\texttt{╟───┼─┘◌║ 6} \\
\texttt{║D◌T│P◌◌║ 7} \\
\texttt{╚═══╧═══╝} \\
\texttt{```} \\
\texttt{Pašlaik jūs spēlējat kārtu 4 no maksimāli 28.} \\
\texttt{Jūs kopā esat uzkrājuši 1 no 16 sodiem.} \\
\texttt{Otrs spēlētājs pārvietoja objektu savā režģī.} \\
\\ 
\\ 
\texttt{Kāda ir jūsu nākamā komanda?} \\
            }
        }
    }
    & \\ \\

    \theutterance \stepcounter{utterance}  
    & & & \multicolumn{4}{p{0.6\linewidth}}{
        \cellcolor[rgb]{0.9,0.9,0.9}{
            \makecell[{{p{\linewidth}}}]{
                \texttt{\tiny{[GM$\langle$P2]}}
                \texttt{GĀJIENS: P, (1,1)} \\
            }
        }
    }
     \\ \\

    \theutterance \stepcounter{utterance}  
    & & & \multicolumn{2}{p{0.3\linewidth}}{
        \cellcolor[rgb]{0.9,0.9,0.9}{
            \makecell[{{p{\linewidth}}}]{
                \texttt{\tiny{[GM$|$GM]}}
                \texttt{GĀJIENS: P, (1,1)} \\
            }
        }
    }
    & & \\ \\

    \theutterance \stepcounter{utterance}  
    & & & \multicolumn{2}{p{0.3\linewidth}}{
        \cellcolor[rgb]{0.9,0.9,0.9}{
            \makecell[{{p{\linewidth}}}]{
                \texttt{\tiny{[GM$|$GM]}}
                \texttt{P moved to (1, 1)} \\
            }
        }
    }
    & & \\ \\

    \theutterance \stepcounter{utterance}  
    & & \multicolumn{4}{p{0.6\linewidth}}{
        \cellcolor[rgb]{0.9,0.9,0.9}{
            \makecell[{{p{\linewidth}}}]{
                \texttt{\tiny{[P1$\langle$GM]}}
                \texttt{Veiksmīgi pārvietots 'P' uz (1,1). Jūsu atjauninātais režģis izskatās šādi:} \\
\texttt{```} \\
\texttt{ 1234567} \\
\texttt{╔═══════╗} \\
\texttt{║P◌◌S◌T◌║ 1} \\
\texttt{║◌◌◌◌◌◌◌║ 2} \\
\texttt{║◌◌AD◌┌─╢ 3} \\
\texttt{║◌◌◌┌─┤◌║ 4} \\
\texttt{║O◌◌│◌├─╢ 5} \\
\texttt{╟───┼─┘◌║ 6} \\
\texttt{║◌M◌│◌◌◌║ 7} \\
\texttt{╚═══╧═══╝} \\
\texttt{```} \\
\texttt{Pašlaik jūs spēlējat kārtu 4 no maksimāli 28.} \\
\texttt{Jūs kopā esat uzkrājuši 1 no 16 sodiem.} \\
\texttt{Otrs spēlētājs pārvietoja objektu savā režģī.} \\
\\ 
\\ 
\texttt{Kāda ir jūsu nākamā komanda?} \\
            }
        }
    }
    & \\ \\

    \theutterance \stepcounter{utterance}  
    & \multicolumn{4}{p{0.6\linewidth}}{
        \cellcolor[rgb]{0.9,0.9,0.9}{
            \makecell[{{p{\linewidth}}}]{
                \texttt{\tiny{[P1$\rangle$GM]}}
                \texttt{GĀJIENS: O, (2,1)} \\
            }
        }
    }
    & & \\ \\

    \theutterance \stepcounter{utterance}  
    & & & \multicolumn{2}{p{0.3\linewidth}}{
        \cellcolor[rgb]{0.9,0.9,0.9}{
            \makecell[{{p{\linewidth}}}]{
                \texttt{\tiny{[GM$|$GM]}}
                \texttt{GĀJIENS: O, (2,1)} \\
            }
        }
    }
    & & \\ \\

    \theutterance \stepcounter{utterance}  
    & & & \multicolumn{2}{p{0.3\linewidth}}{
        \cellcolor[rgb]{0.9,0.9,0.9}{
            \makecell[{{p{\linewidth}}}]{
                \texttt{\tiny{[GM$|$GM]}}
                \texttt{O moved to (2, 1)} \\
            }
        }
    }
    & & \\ \\

    \theutterance \stepcounter{utterance}  
    & & \multicolumn{4}{p{0.6\linewidth}}{
        \cellcolor[rgb]{0.9,0.9,0.9}{
            \makecell[{{p{\linewidth}}}]{
                \texttt{\tiny{[GM$\rangle$P2]}}
                \texttt{Veiksmīgi pārvietots 'P' uz (1,1). Jūsu atjauninātais režģis izskatās šādi:} \\
\texttt{```} \\
\texttt{ 1234567} \\
\texttt{╔═══════╗} \\
\texttt{║P◌◌S◌◌◌║ 1} \\
\texttt{║◌◌◌◌O◌◌║ 2} \\
\texttt{║◌◌◌M◌┌─╢ 3} \\
\texttt{║◌◌A┌─┤◌║ 4} \\
\texttt{║◌◌◌│◌├─╢ 5} \\
\texttt{╟───┼─┘◌║ 6} \\
\texttt{║D◌T│◌◌◌║ 7} \\
\texttt{╚═══╧═══╝} \\
\texttt{```} \\
\texttt{Pašlaik jūs spēlējat kārtu 5 no maksimāli 28.} \\
\texttt{Jūs kopā esat uzkrājuši 1 no 16 sodiem.} \\
\texttt{Otrs spēlētājs pārvietoja objektu savā režģī.} \\
\\ 
\\ 
\texttt{Kāda ir jūsu nākamā komanda?} \\
            }
        }
    }
    & \\ \\

    \theutterance \stepcounter{utterance}  
    & & & \multicolumn{4}{p{0.6\linewidth}}{
        \cellcolor[rgb]{0.9,0.9,0.9}{
            \makecell[{{p{\linewidth}}}]{
                \texttt{\tiny{[GM$\langle$P2]}}
                \texttt{GĀJIENS: O, (2,1)} \\
            }
        }
    }
     \\ \\

    \theutterance \stepcounter{utterance}  
    & & & \multicolumn{2}{p{0.3\linewidth}}{
        \cellcolor[rgb]{0.9,0.9,0.9}{
            \makecell[{{p{\linewidth}}}]{
                \texttt{\tiny{[GM$|$GM]}}
                \texttt{GĀJIENS: O, (2,1)} \\
            }
        }
    }
    & & \\ \\

    \theutterance \stepcounter{utterance}  
    & & & \multicolumn{2}{p{0.3\linewidth}}{
        \cellcolor[rgb]{0.9,0.9,0.9}{
            \makecell[{{p{\linewidth}}}]{
                \texttt{\tiny{[GM$|$GM]}}
                \texttt{O moved to (2, 1)} \\
            }
        }
    }
    & & \\ \\

    \theutterance \stepcounter{utterance}  
    & & \multicolumn{4}{p{0.6\linewidth}}{
        \cellcolor[rgb]{0.9,0.9,0.9}{
            \makecell[{{p{\linewidth}}}]{
                \texttt{\tiny{[P1$\langle$GM]}}
                \texttt{Veiksmīgi pārvietots 'O' uz (2,1). Jūsu atjauninātais režģis izskatās šādi:} \\
\texttt{```} \\
\texttt{ 1234567} \\
\texttt{╔═══════╗} \\
\texttt{║PO◌S◌T◌║ 1} \\
\texttt{║◌◌◌◌◌◌◌║ 2} \\
\texttt{║◌◌AD◌┌─╢ 3} \\
\texttt{║◌◌◌┌─┤◌║ 4} \\
\texttt{║◌◌◌│◌├─╢ 5} \\
\texttt{╟───┼─┘◌║ 6} \\
\texttt{║◌M◌│◌◌◌║ 7} \\
\texttt{╚═══╧═══╝} \\
\texttt{```} \\
\texttt{Pašlaik jūs spēlējat kārtu 5 no maksimāli 28.} \\
\texttt{Jūs kopā esat uzkrājuši 1 no 16 sodiem.} \\
\texttt{Otrs spēlētājs pārvietoja objektu savā režģī.} \\
\\ 
\\ 
\texttt{Kāda ir jūsu nākamā komanda?} \\
            }
        }
    }
    & \\ \\

    \theutterance \stepcounter{utterance}  
    & \multicolumn{4}{p{0.6\linewidth}}{
        \cellcolor[rgb]{0.9,0.9,0.9}{
            \makecell[{{p{\linewidth}}}]{
                \texttt{\tiny{[P1$\rangle$GM]}}
                \texttt{GĀJIENS: T, (3,1)} \\
            }
        }
    }
    & & \\ \\

    \theutterance \stepcounter{utterance}  
    & & & \multicolumn{2}{p{0.3\linewidth}}{
        \cellcolor[rgb]{0.9,0.9,0.9}{
            \makecell[{{p{\linewidth}}}]{
                \texttt{\tiny{[GM$|$GM]}}
                \texttt{GĀJIENS: T, (3,1)} \\
            }
        }
    }
    & & \\ \\

    \theutterance \stepcounter{utterance}  
    & & & \multicolumn{2}{p{0.3\linewidth}}{
        \cellcolor[rgb]{0.9,0.9,0.9}{
            \makecell[{{p{\linewidth}}}]{
                \texttt{\tiny{[GM$|$GM]}}
                \texttt{T moved to (3, 1)} \\
            }
        }
    }
    & & \\ \\

    \theutterance \stepcounter{utterance}  
    & & \multicolumn{4}{p{0.6\linewidth}}{
        \cellcolor[rgb]{0.9,0.9,0.9}{
            \makecell[{{p{\linewidth}}}]{
                \texttt{\tiny{[GM$\rangle$P2]}}
                \texttt{Veiksmīgi pārvietots 'O' uz (2,1). Jūsu atjauninātais režģis izskatās šādi:} \\
\texttt{```} \\
\texttt{ 1234567} \\
\texttt{╔═══════╗} \\
\texttt{║PO◌S◌◌◌║ 1} \\
\texttt{║◌◌◌◌◌◌◌║ 2} \\
\texttt{║◌◌◌M◌┌─╢ 3} \\
\texttt{║◌◌A┌─┤◌║ 4} \\
\texttt{║◌◌◌│◌├─╢ 5} \\
\texttt{╟───┼─┘◌║ 6} \\
\texttt{║D◌T│◌◌◌║ 7} \\
\texttt{╚═══╧═══╝} \\
\texttt{```} \\
\texttt{Pašlaik jūs spēlējat kārtu 6 no maksimāli 28.} \\
\texttt{Jūs kopā esat uzkrājuši 1 no 16 sodiem.} \\
\texttt{Otrs spēlētājs pārvietoja objektu savā režģī.} \\
\\ 
\\ 
\texttt{Kāda ir jūsu nākamā komanda?} \\
            }
        }
    }
    & \\ \\

    \theutterance \stepcounter{utterance}  
    & & & \multicolumn{4}{p{0.6\linewidth}}{
        \cellcolor[rgb]{0.9,0.9,0.9}{
            \makecell[{{p{\linewidth}}}]{
                \texttt{\tiny{[GM$\langle$P2]}}
                \texttt{GĀJIENS: T, (3,1)} \\
            }
        }
    }
     \\ \\

    \theutterance \stepcounter{utterance}  
    & & & \multicolumn{2}{p{0.3\linewidth}}{
        \cellcolor[rgb]{0.9,0.9,0.9}{
            \makecell[{{p{\linewidth}}}]{
                \texttt{\tiny{[GM$|$GM]}}
                \texttt{GĀJIENS: T, (3,1)} \\
            }
        }
    }
    & & \\ \\

    \theutterance \stepcounter{utterance}  
    & & & \multicolumn{2}{p{0.3\linewidth}}{
        \cellcolor[rgb]{0.9,0.9,0.9}{
            \makecell[{{p{\linewidth}}}]{
                \texttt{\tiny{[GM$|$GM]}}
                \texttt{T moved to (3, 1)} \\
            }
        }
    }
    & & \\ \\

    \theutterance \stepcounter{utterance}  
    & & \multicolumn{4}{p{0.6\linewidth}}{
        \cellcolor[rgb]{0.9,0.9,0.9}{
            \makecell[{{p{\linewidth}}}]{
                \texttt{\tiny{[P1$\langle$GM]}}
                \texttt{Veiksmīgi pārvietots 'T' uz (3,1). Jūsu atjauninātais režģis izskatās šādi:} \\
\texttt{```} \\
\texttt{ 1234567} \\
\texttt{╔═══════╗} \\
\texttt{║POTS◌◌◌║ 1} \\
\texttt{║◌◌◌◌◌◌◌║ 2} \\
\texttt{║◌◌AD◌┌─╢ 3} \\
\texttt{║◌◌◌┌─┤◌║ 4} \\
\texttt{║◌◌◌│◌├─╢ 5} \\
\texttt{╟───┼─┘◌║ 6} \\
\texttt{║◌M◌│◌◌◌║ 7} \\
\texttt{╚═══╧═══╝} \\
\texttt{```} \\
\texttt{Pašlaik jūs spēlējat kārtu 6 no maksimāli 28.} \\
\texttt{Jūs kopā esat uzkrājuši 1 no 16 sodiem.} \\
\texttt{Otrs spēlētājs pārvietoja objektu savā režģī.} \\
\\ 
\\ 
\texttt{Kāda ir jūsu nākamā komanda?} \\
            }
        }
    }
    & \\ \\

    \theutterance \stepcounter{utterance}  
    & \multicolumn{4}{p{0.6\linewidth}}{
        \cellcolor[rgb]{0.9,0.9,0.9}{
            \makecell[{{p{\linewidth}}}]{
                \texttt{\tiny{[P1$\rangle$GM]}}
                \texttt{GĀJIENS: D, (5,1)} \\
            }
        }
    }
    & & \\ \\

    \theutterance \stepcounter{utterance}  
    & & & \multicolumn{2}{p{0.3\linewidth}}{
        \cellcolor[rgb]{0.9,0.9,0.9}{
            \makecell[{{p{\linewidth}}}]{
                \texttt{\tiny{[GM$|$GM]}}
                \texttt{GĀJIENS: D, (5,1)} \\
            }
        }
    }
    & & \\ \\

    \theutterance \stepcounter{utterance}  
    & & & \multicolumn{2}{p{0.3\linewidth}}{
        \cellcolor[rgb]{0.9,0.9,0.9}{
            \makecell[{{p{\linewidth}}}]{
                \texttt{\tiny{[GM$|$GM]}}
                \texttt{D moved to (5, 1)} \\
            }
        }
    }
    & & \\ \\

    \theutterance \stepcounter{utterance}  
    & & \multicolumn{4}{p{0.6\linewidth}}{
        \cellcolor[rgb]{0.9,0.9,0.9}{
            \makecell[{{p{\linewidth}}}]{
                \texttt{\tiny{[GM$\rangle$P2]}}
                \texttt{Veiksmīgi pārvietots 'T' uz (3,1). Jūsu atjauninātais režģis izskatās šādi:} \\
\texttt{```} \\
\texttt{ 1234567} \\
\texttt{╔═══════╗} \\
\texttt{║POTS◌◌◌║ 1} \\
\texttt{║◌◌◌◌◌◌◌║ 2} \\
\texttt{║◌◌◌M◌┌─╢ 3} \\
\texttt{║◌◌A┌─┤◌║ 4} \\
\texttt{║◌◌◌│◌├─╢ 5} \\
\texttt{╟───┼─┘◌║ 6} \\
\texttt{║D◌◌│◌◌◌║ 7} \\
\texttt{╚═══╧═══╝} \\
\texttt{```} \\
\texttt{Pašlaik jūs spēlējat kārtu 7 no maksimāli 28.} \\
\texttt{Jūs kopā esat uzkrājuši 1 no 16 sodiem.} \\
\texttt{Otrs spēlētājs pārvietoja objektu savā režģī.} \\
\\ 
\\ 
\texttt{Kāda ir jūsu nākamā komanda?} \\
            }
        }
    }
    & \\ \\

    \theutterance \stepcounter{utterance}  
    & & & \multicolumn{4}{p{0.6\linewidth}}{
        \cellcolor[rgb]{0.9,0.9,0.9}{
            \makecell[{{p{\linewidth}}}]{
                \texttt{\tiny{[GM$\langle$P2]}}
                \texttt{GĀJIENS: D, (5,1)} \\
            }
        }
    }
     \\ \\

    \theutterance \stepcounter{utterance}  
    & & & \multicolumn{2}{p{0.3\linewidth}}{
        \cellcolor[rgb]{0.9,0.9,0.9}{
            \makecell[{{p{\linewidth}}}]{
                \texttt{\tiny{[GM$|$GM]}}
                \texttt{GĀJIENS: D, (5,1)} \\
            }
        }
    }
    & & \\ \\

    \theutterance \stepcounter{utterance}  
    & & & \multicolumn{2}{p{0.3\linewidth}}{
        \cellcolor[rgb]{0.9,0.9,0.9}{
            \makecell[{{p{\linewidth}}}]{
                \texttt{\tiny{[GM$|$GM]}}
                \texttt{D moved to (5, 1)} \\
            }
        }
    }
    & & \\ \\

    \theutterance \stepcounter{utterance}  
    & & \multicolumn{4}{p{0.6\linewidth}}{
        \cellcolor[rgb]{0.9,0.9,0.9}{
            \makecell[{{p{\linewidth}}}]{
                \texttt{\tiny{[P1$\langle$GM]}}
                \texttt{Veiksmīgi pārvietots 'D' uz (5,1). Jūsu atjauninātais režģis izskatās šādi:} \\
\texttt{```} \\
\texttt{ 1234567} \\
\texttt{╔═══════╗} \\
\texttt{║POTSD◌◌║ 1} \\
\texttt{║◌◌◌◌◌◌◌║ 2} \\
\texttt{║◌◌A◌◌┌─╢ 3} \\
\texttt{║◌◌◌┌─┤◌║ 4} \\
\texttt{║◌◌◌│◌├─╢ 5} \\
\texttt{╟───┼─┘◌║ 6} \\
\texttt{║◌M◌│◌◌◌║ 7} \\
\texttt{╚═══╧═══╝} \\
\texttt{```} \\
\texttt{Pašlaik jūs spēlējat kārtu 7 no maksimāli 28.} \\
\texttt{Jūs kopā esat uzkrājuši 1 no 16 sodiem.} \\
\texttt{Otrs spēlētājs pārvietoja objektu savā režģī.} \\
\\ 
\\ 
\texttt{Kāda ir jūsu nākamā komanda?} \\
            }
        }
    }
    & \\ \\

    \theutterance \stepcounter{utterance}  
    & \multicolumn{4}{p{0.6\linewidth}}{
        \cellcolor[rgb]{0.9,0.9,0.9}{
            \makecell[{{p{\linewidth}}}]{
                \texttt{\tiny{[P1$\rangle$GM]}}
                \texttt{GĀJIENS: A, (6,1)} \\
            }
        }
    }
    & & \\ \\

    \theutterance \stepcounter{utterance}  
    & & & \multicolumn{2}{p{0.3\linewidth}}{
        \cellcolor[rgb]{0.9,0.9,0.9}{
            \makecell[{{p{\linewidth}}}]{
                \texttt{\tiny{[GM$|$GM]}}
                \texttt{GĀJIENS: A, (6,1)} \\
            }
        }
    }
    & & \\ \\

    \theutterance \stepcounter{utterance}  
    & & & \multicolumn{2}{p{0.3\linewidth}}{
        \cellcolor[rgb]{0.9,0.9,0.9}{
            \makecell[{{p{\linewidth}}}]{
                \texttt{\tiny{[GM$|$GM]}}
                \texttt{A moved to (6, 1)} \\
            }
        }
    }
    & & \\ \\

    \theutterance \stepcounter{utterance}  
    & & \multicolumn{4}{p{0.6\linewidth}}{
        \cellcolor[rgb]{0.9,0.9,0.9}{
            \makecell[{{p{\linewidth}}}]{
                \texttt{\tiny{[GM$\rangle$P2]}}
                \texttt{Veiksmīgi pārvietots 'D' uz (5,1). Jūsu atjauninātais režģis izskatās šādi:} \\
\texttt{```} \\
\texttt{ 1234567} \\
\texttt{╔═══════╗} \\
\texttt{║POTSD◌◌║ 1} \\
\texttt{║◌◌◌◌◌◌◌║ 2} \\
\texttt{║◌◌◌M◌┌─╢ 3} \\
\texttt{║◌◌A┌─┤◌║ 4} \\
\texttt{║◌◌◌│◌├─╢ 5} \\
\texttt{╟───┼─┘◌║ 6} \\
\texttt{║◌◌◌│◌◌◌║ 7} \\
\texttt{╚═══╧═══╝} \\
\texttt{```} \\
\texttt{Pašlaik jūs spēlējat kārtu 8 no maksimāli 28.} \\
\texttt{Jūs kopā esat uzkrājuši 1 no 16 sodiem.} \\
\texttt{Otrs spēlētājs pārvietoja objektu savā režģī.} \\
\\ 
\\ 
\texttt{Kāda ir jūsu nākamā komanda?} \\
            }
        }
    }
    & \\ \\

    \theutterance \stepcounter{utterance}  
    & & & \multicolumn{4}{p{0.6\linewidth}}{
        \cellcolor[rgb]{0.9,0.9,0.9}{
            \makecell[{{p{\linewidth}}}]{
                \texttt{\tiny{[GM$\langle$P2]}}
                \texttt{GĀJIENS: A, (6,1)} \\
            }
        }
    }
     \\ \\

    \theutterance \stepcounter{utterance}  
    & & & \multicolumn{2}{p{0.3\linewidth}}{
        \cellcolor[rgb]{0.9,0.9,0.9}{
            \makecell[{{p{\linewidth}}}]{
                \texttt{\tiny{[GM$|$GM]}}
                \texttt{GĀJIENS: A, (6,1)} \\
            }
        }
    }
    & & \\ \\

    \theutterance \stepcounter{utterance}  
    & & & \multicolumn{2}{p{0.3\linewidth}}{
        \cellcolor[rgb]{0.9,0.9,0.9}{
            \makecell[{{p{\linewidth}}}]{
                \texttt{\tiny{[GM$|$GM]}}
                \texttt{A moved to (6, 1)} \\
            }
        }
    }
    & & \\ \\

    \theutterance \stepcounter{utterance}  
    & & \multicolumn{4}{p{0.6\linewidth}}{
        \cellcolor[rgb]{0.9,0.9,0.9}{
            \makecell[{{p{\linewidth}}}]{
                \texttt{\tiny{[P1$\langle$GM]}}
                \texttt{Veiksmīgi pārvietots 'A' uz (6,1). Jūsu atjauninātais režģis izskatās šādi:} \\
\texttt{```} \\
\texttt{ 1234567} \\
\texttt{╔═══════╗} \\
\texttt{║POTSDA◌║ 1} \\
\texttt{║◌◌◌◌◌◌◌║ 2} \\
\texttt{║◌◌◌◌◌┌─╢ 3} \\
\texttt{║◌◌◌┌─┤◌║ 4} \\
\texttt{║◌◌◌│◌├─╢ 5} \\
\texttt{╟───┼─┘◌║ 6} \\
\texttt{║◌M◌│◌◌◌║ 7} \\
\texttt{╚═══╧═══╝} \\
\texttt{```} \\
\texttt{Pašlaik jūs spēlējat kārtu 8 no maksimāli 28.} \\
\texttt{Jūs kopā esat uzkrājuši 1 no 16 sodiem.} \\
\texttt{Otrs spēlētājs pārvietoja objektu savā režģī.} \\
\\ 
\\ 
\texttt{Kāda ir jūsu nākamā komanda?} \\
            }
        }
    }
    & \\ \\

    \theutterance \stepcounter{utterance}  
    & \multicolumn{4}{p{0.6\linewidth}}{
        \cellcolor[rgb]{0.9,0.9,0.9}{
            \makecell[{{p{\linewidth}}}]{
                \texttt{\tiny{[P1$\rangle$GM]}}
                \texttt{GĀJIENS: M, (1,2)} \\
            }
        }
    }
    & & \\ \\

    \theutterance \stepcounter{utterance}  
    & & & \multicolumn{2}{p{0.3\linewidth}}{
        \cellcolor[rgb]{0.9,0.9,0.9}{
            \makecell[{{p{\linewidth}}}]{
                \texttt{\tiny{[GM$|$GM]}}
                \texttt{GĀJIENS: M, (1,2)} \\
            }
        }
    }
    & & \\ \\

    \theutterance \stepcounter{utterance}  
    & & & \multicolumn{2}{p{0.3\linewidth}}{
        \cellcolor[rgb]{0.9,0.9,0.9}{
            \makecell[{{p{\linewidth}}}]{
                \texttt{\tiny{[GM$|$GM]}}
                \texttt{M moved to (1, 2)} \\
            }
        }
    }
    & & \\ \\

    \theutterance \stepcounter{utterance}  
    & & \multicolumn{4}{p{0.6\linewidth}}{
        \cellcolor[rgb]{0.9,0.9,0.9}{
            \makecell[{{p{\linewidth}}}]{
                \texttt{\tiny{[GM$\rangle$P2]}}
                \texttt{Veiksmīgi pārvietots 'A' uz (6,1). Jūsu atjauninātais režģis izskatās šādi:} \\
\texttt{```} \\
\texttt{ 1234567} \\
\texttt{╔═══════╗} \\
\texttt{║POTSDA◌║ 1} \\
\texttt{║◌◌◌◌◌◌◌║ 2} \\
\texttt{║◌◌◌M◌┌─╢ 3} \\
\texttt{║◌◌◌┌─┤◌║ 4} \\
\texttt{║◌◌◌│◌├─╢ 5} \\
\texttt{╟───┼─┘◌║ 6} \\
\texttt{║◌◌◌│◌◌◌║ 7} \\
\texttt{╚═══╧═══╝} \\
\texttt{```} \\
\texttt{Pašlaik jūs spēlējat kārtu 9 no maksimāli 28.} \\
\texttt{Jūs kopā esat uzkrājuši 1 no 16 sodiem.} \\
\texttt{Otrs spēlētājs pārvietoja objektu savā režģī.} \\
\\ 
\\ 
\texttt{Kāda ir jūsu nākamā komanda?} \\
            }
        }
    }
    & \\ \\

    \theutterance \stepcounter{utterance}  
    & & & \multicolumn{4}{p{0.6\linewidth}}{
        \cellcolor[rgb]{0.9,0.9,0.9}{
            \makecell[{{p{\linewidth}}}]{
                \texttt{\tiny{[GM$\langle$P2]}}
                \texttt{GĀJIENS: M, (1,2)} \\
            }
        }
    }
     \\ \\

    \theutterance \stepcounter{utterance}  
    & & & \multicolumn{2}{p{0.3\linewidth}}{
        \cellcolor[rgb]{0.9,0.9,0.9}{
            \makecell[{{p{\linewidth}}}]{
                \texttt{\tiny{[GM$|$GM]}}
                \texttt{GĀJIENS: M, (1,2)} \\
            }
        }
    }
    & & \\ \\

    \theutterance \stepcounter{utterance}  
    & & & \multicolumn{2}{p{0.3\linewidth}}{
        \cellcolor[rgb]{0.9,0.9,0.9}{
            \makecell[{{p{\linewidth}}}]{
                \texttt{\tiny{[GM$|$GM]}}
                \texttt{M moved to (1, 2)} \\
            }
        }
    }
    & & \\ \\

    \theutterance \stepcounter{utterance}  
    & & \multicolumn{4}{p{0.6\linewidth}}{
        \cellcolor[rgb]{0.9,0.9,0.9}{
            \makecell[{{p{\linewidth}}}]{
                \texttt{\tiny{[P1$\langle$GM]}}
                \texttt{Veiksmīgi pārvietots 'M' uz (1,2). Jūsu atjauninātais režģis izskatās šādi:} \\
\texttt{```} \\
\texttt{ 1234567} \\
\texttt{╔═══════╗} \\
\texttt{║POTSDA◌║ 1} \\
\texttt{║M◌◌◌◌◌◌║ 2} \\
\texttt{║◌◌◌◌◌┌─╢ 3} \\
\texttt{║◌◌◌┌─┤◌║ 4} \\
\texttt{║◌◌◌│◌├─╢ 5} \\
\texttt{╟───┼─┘◌║ 6} \\
\texttt{║◌◌◌│◌◌◌║ 7} \\
\texttt{╚═══╧═══╝} \\
\texttt{```} \\
\texttt{Pašlaik jūs spēlējat kārtu 9 no maksimāli 28.} \\
\texttt{Jūs kopā esat uzkrājuši 1 no 16 sodiem.} \\
\texttt{Otrs spēlētājs pārvietoja objektu savā režģī.} \\
\\ 
\\ 
\texttt{Kāda ir jūsu nākamā komanda?} \\
            }
        }
    }
    & \\ \\

    \theutterance \stepcounter{utterance}  
    & \multicolumn{4}{p{0.6\linewidth}}{
        \cellcolor[rgb]{0.9,0.9,0.9}{
            \makecell[{{p{\linewidth}}}]{
                \texttt{\tiny{[P1$\rangle$GM]}}
                \texttt{SAKI: pabeigts?} \\
            }
        }
    }
    & & \\ \\

    \theutterance \stepcounter{utterance}  
    & & & \multicolumn{2}{p{0.3\linewidth}}{
        \cellcolor[rgb]{0.9,0.9,0.9}{
            \makecell[{{p{\linewidth}}}]{
                \texttt{\tiny{[GM$|$GM]}}
                \texttt{SAKI: pabeigts?} \\
            }
        }
    }
    & & \\ \\

    \theutterance \stepcounter{utterance}  
    & & \multicolumn{4}{p{0.6\linewidth}}{
        \cellcolor[rgb]{0.9,0.9,0.9}{
            \makecell[{{p{\linewidth}}}]{
                \texttt{\tiny{[GM$\rangle$P2]}}
                \texttt{Veiksmīgi pārvietots 'M' uz (1,2). Jūsu atjauninātais režģis izskatās šādi:} \\
\texttt{```} \\
\texttt{ 1234567} \\
\texttt{╔═══════╗} \\
\texttt{║POTSDA◌║ 1} \\
\texttt{║M◌◌◌◌◌◌║ 2} \\
\texttt{║◌◌◌◌◌┌─╢ 3} \\
\texttt{║◌◌◌┌─┤◌║ 4} \\
\texttt{║◌◌◌│◌├─╢ 5} \\
\texttt{╟───┼─┘◌║ 6} \\
\texttt{║◌◌◌│◌◌◌║ 7} \\
\texttt{╚═══╧═══╝} \\
\texttt{```} \\
\texttt{Pašlaik jūs spēlējat kārtu 10 no maksimāli 28.} \\
\texttt{Jūs kopā esat uzkrājuši 1 no 16 sodiem.} \\
\texttt{Otrs spēlētājs nosūtīja šo ziņojumu:} \\
\\ 
\\ 
\texttt{"pabeigts?"} \\
\\ 
\\ 
\texttt{Kāda ir jūsu nākamā komanda?} \\
            }
        }
    }
    & \\ \\

    \theutterance \stepcounter{utterance}  
    & & & \multicolumn{4}{p{0.6\linewidth}}{
        \cellcolor[rgb]{0.9,0.9,0.9}{
            \makecell[{{p{\linewidth}}}]{
                \texttt{\tiny{[GM$\langle$P2]}}
                \texttt{SAKI: pabeigts!} \\
            }
        }
    }
     \\ \\

    \theutterance \stepcounter{utterance}  
    & & & \multicolumn{2}{p{0.3\linewidth}}{
        \cellcolor[rgb]{0.9,0.9,0.9}{
            \makecell[{{p{\linewidth}}}]{
                \texttt{\tiny{[GM$|$GM]}}
                \texttt{SAKI: pabeigts!} \\
            }
        }
    }
    & & \\ \\

    \theutterance \stepcounter{utterance}  
    & & & \multicolumn{2}{p{0.3\linewidth}}{
        \cellcolor[rgb]{0.9,0.9,0.9}{
            \makecell[{{p{\linewidth}}}]{
                \texttt{\tiny{[GM$|$GM]}}
                \texttt{true} \\
            }
        }
    }
    & & \\ \\

    \theutterance \stepcounter{utterance}  
    & & & \multicolumn{2}{p{0.3\linewidth}}{
        \cellcolor[rgb]{0.9,0.9,0.9}{
            \makecell[{{p{\linewidth}}}]{
                \texttt{\tiny{[GM$|$GM]}}
                \texttt{Initial states:} \\
\texttt{```} \\
\texttt{Player 1:} \\
\texttt{ 1234567} \\
\texttt{╔═══════╗} \\
\texttt{║◌◌◌◌◌T◌║ 1} \\
\texttt{║P◌◌◌◌◌◌║ 2} \\
\texttt{║◌◌AD◌┌─╢ 3} \\
\texttt{║◌◌◌┌─┤S║ 4} \\
\texttt{║O◌◌│◌├─╢ 5} \\
\texttt{╟───┼─┘◌║ 6} \\
\texttt{║◌M◌│◌◌◌║ 7} \\
\texttt{╚═══╧═══╝} \\
\\ 
\\ 
\texttt{Player 2:} \\
\texttt{ 1234567} \\
\texttt{╔═══════╗} \\
\texttt{║◌◌◌◌◌S◌║ 1} \\
\texttt{║◌◌◌◌O◌◌║ 2} \\
\texttt{║◌◌◌M◌┌─╢ 3} \\
\texttt{║◌◌A┌─┤◌║ 4} \\
\texttt{║◌◌◌│◌├─╢ 5} \\
\texttt{╟───┼─┘◌║ 6} \\
\texttt{║D◌T│P◌◌║ 7} \\
\texttt{╚═══╧═══╝} \\
\\ 
\texttt{```} \\
            }
        }
    }
    & & \\ \\

    \theutterance \stepcounter{utterance}  
    & & & \multicolumn{2}{p{0.3\linewidth}}{
        \cellcolor[rgb]{0.9,0.9,0.9}{
            \makecell[{{p{\linewidth}}}]{
                \texttt{\tiny{[GM$|$GM]}}
                \texttt{End states:} \\
\texttt{```} \\
\texttt{Player 1:} \\
\texttt{ 1234567} \\
\texttt{╔═══════╗} \\
\texttt{║POTSDA◌║ 1} \\
\texttt{║M◌◌◌◌◌◌║ 2} \\
\texttt{║◌◌◌◌◌┌─╢ 3} \\
\texttt{║◌◌◌┌─┤◌║ 4} \\
\texttt{║◌◌◌│◌├─╢ 5} \\
\texttt{╟───┼─┘◌║ 6} \\
\texttt{║◌◌◌│◌◌◌║ 7} \\
\texttt{╚═══╧═══╝} \\
\\ 
\\ 
\texttt{Player 2:} \\
\texttt{ 1234567} \\
\texttt{╔═══════╗} \\
\texttt{║POTSDA◌║ 1} \\
\texttt{║M◌◌◌◌◌◌║ 2} \\
\texttt{║◌◌◌◌◌┌─╢ 3} \\
\texttt{║◌◌◌┌─┤◌║ 4} \\
\texttt{║◌◌◌│◌├─╢ 5} \\
\texttt{╟───┼─┘◌║ 6} \\
\texttt{║◌◌◌│◌◌◌║ 7} \\
\texttt{╚═══╧═══╝} \\
\\ 
\texttt{```} \\
            }
        }
    }
    & & \\ \\

    \theutterance \stepcounter{utterance}  
    & & & \multicolumn{2}{p{0.3\linewidth}}{
        \cellcolor[rgb]{0.9,0.9,0.9}{
            \makecell[{{p{\linewidth}}}]{
                \texttt{\tiny{[GM$|$GM]}}
                \texttt{* success: True} \\
\texttt{* lose: False} \\
\texttt{* aborted: False} \\
\texttt{{-}{-}{-}{-}{-}{-}{-}} \\
\texttt{* Shifts: 6.00} \\
\texttt{* Max Shifts: 12.00} \\
\texttt{* Min Shifts: 6.00} \\
\texttt{* End Distance Sum: 0.00} \\
\texttt{* Init Distance Sum: 31.75} \\
\texttt{* Expected Distance Sum: 29.33} \\
\texttt{* Penalties: 1.00} \\
\texttt{* Max Penalties: 16.00} \\
\texttt{* Rounds: 10.00} \\
\texttt{* Max Rounds: 28.00} \\
\texttt{* Object Count: 7.00} \\
            }
        }
    }
    & & \\ \\

\end{supertabular}
}

\end{document}
